\documentclass[11pt]{book}

\usepackage{graphicx}
\usepackage{eso-pic}
\usepackage{tcolorbox}
\usepackage{polyglossia}
\setdefaultlanguage[indentfirst=false,forceheadingpunctuation=false]{russian}
\setotherlanguages{english}

\RequirePackage[a4paper, headsep = 0.5 \headsep, left=2.5cm, right=2.1cm, top=2cm, bottom=2.1cm]{geometry}

\setmainfont{Times New Roman}
\newfontfamily\cyrillicfont{Times New Roman}[Script=Cyrillic, Ligatures=TeX]

\usepackage{enumitem}
\setlist[itemize, 1]{label =\raisebox{-0.3\height}{\scalebox{1.8}{\textbullet}}}

\title{Естественные основания}
\author{Сьюзен Хёрли}
\date{1989}
\begin{document}

\AddToShipoutPictureBG*{\includegraphics[width=\paperwidth]{title.png}}

\null
    \thispagestyle{empty}%
    \addtocounter{page}{-1}%
    \newpage

\tableofcontents


\chapter{Введение}

Эта работа посвящена рациональности решений и действий людей и иллюстрации преемственности философии разума, с одной стороны, и этики и юриспруденции, с другой. Она не о морали, и, более того, отвергает традиционную ее и личных интересов концепцию как существенно раздельных. Речь, скорее, пойдет об этике или рациональности со стороны самоинтерпретирующих и сапоопределяющихся агентов. Главный тезис --- аргументы из философии сознания могут подорвать распространенные субъективистские позиции в этике и соответствующие позиции в политико-экономической теории. Книга эта во многом вдохновлена работами Витгенштейна и Дэвидсона, и связывает их аргументы касательно интерпретации с теорией принятия решений, теорией социального выбора и теорией вынесения судебных решений. На протяжении всей книги будут рассматриваться сходства и несоответствия межличностных и внутриличностных отношений, межличностных и внутриличностных конфликтов, и тема эта завершится в последней главе рассмотрением тесной связи между ценностями личной автономии и демократии.

Книга состоит из четырех частей. Часть I посвящена соотнесению разума и ценности и содержит главы об объективности, разногласиях, предпочтениях, интерпретации и субъективизме. Представленные в ней аргументы касательно интерпретации заимствованы из философии сознания и сосредоточены на концепции предпочтения из теории принятия решений. Теория принятия решений допускает ошибки и разногласия по поводу того, что следует делать. Для этого должны быть ограничения на интерпретации действий в терминах убеждений и предпочтений. Ограничения нужны как существенные, так и формальные, в том числе оценочные. Будут разграничены вопросы об объективности, реализме и когнитивизме, а также подвергнуто критике субъективистское утверждение о детерминированности ценностей концептуально первичными и независимыми предпочтениями, и будет показано, что некоторые теории предпочтений высшего порядка и межличностного сравнения зависят от теорий человеческой природы, носящий частично оценочный характер, и что разногласия касательно таких теорий важно формулировать. Попытки субъективистов использовать предпочтения индивидов в качестве беспроблемных строительных блоков для своей теории подрываются конститутивной ролью ценностей в интерпретации. Объективистская концепция предпочтений как ''обремененных'' ценностями будет отделена от платонизма и соотнесена с концепцией разума как обремененного миром, возникающей из современной критики психологического индивидуализма.

Часть II посвящена конфликтующим ценностям и основаниям и содержит главы о конфликте, акрасии и когнитивизме. В ней обосновывается необходимость рассмотрения конфликта оснований к действиям как не просто кажущегося, и доказательная модель конфликта различается с подсистемной, говорящей об аналогичности внутриличностных конфликтов межличностным. Также рассматриваются связи доказательной и подсистемной моделей конфликта с атрибуцией установок акратической структурой, и подчеркивается рациональность внутриличностной и межличностной координации и неадекватность индивидуализма в отношении агентности в его понимании единицы агентности как фиксированной. Люди понимаются как самоинтерпретирующие и самоопределяющиеся животные с нефиксированными и частично ими самими создаваемыми взглядами, характерами и идентичностями; теория рациональности личности тогда и есть этика. Утверждается, что конфликты ценностей в рамках подсистемного подхода совместимы с этическим когнитивизмом и выдвигается аргумент Фреге против нон-когнитивизма.

Часть III посвящена рациональности в условиях противоречивых оснований и содержит главы ''Теория'', ''Делиберация'', ''Когерентность'' и ''Соизмеримость''. Делиберация касательно того, что, учитывая все обстоятельства, следует предпринять, характеризуется как поиск когерентного набора отношений между соответствующими кофликтующими причинами. В этическом контексте это можно понимать как своего рода нашу самоинтерпретацию и самоопределение. Теория, демонстрирующая искомую когерентность, может быть представлена ее функцией из упорядочений альтернатив под конфликтующим основаниям в суждения о том, что, учитывая все обстоятельства, следует сделать. Впрочем, теория социального выбора при определенных ограничениях демонстрирует отсутствие функций социального благосостояния, которые рационально разрешали бы конфликты предпочтений различных агентов. Аналогия между внутриличностными и межличностными конфликтами наводит на вопрос применимости аналогов определенных результатов невозможности в теории социального выбора к поиску когерентности в условиях конфликтующих оснований. Чтобы понять, почему ответ на этот вопрос отрицательный, будет исследована модальная структура делиберации: она включает обращение к контрфактическим, но не к контроценочным возможностям, а поиск когерентности ограничен как конкретными ценностями, так и формальным требованием супервентности оценочных концепций. Эти особенности делиберации опровергают скептическое предположение указывания аналогов некоторых теорем социального выбора на невозможность когерентности, и, значит, того, что следует делать, учитывая все обстоятельства. Более того --- когерентность возможна и без соизмеримости.

Часть IV посвящена знанию о том, что следует делать, и взаимосвязи личной автономии и демократии с этим знанием. В этой части обсуждение реализма и критика теории ошибок касательно ценностей приводят к рассмотрению эпистемологического скептицизма и защите возможности отслеживания истины о том, что следует делать. Скептик пытается развенчать оценочные убеждения, показав необходимость их придерживания даже при их ложности. Понимание модальной структуры делиберации и роли супервентности, достигнутое в Части III, позволяет сопротивляться таким не столь уж проницательным попыткам. Есть, конечно, понимания оценочных убеждений, не ссылающиеся на ценности, которые эти убеждения действительно опровергают, но не все понимания таковы: демократическое разделение эпистемического труда может обеспечить средство избежать влияния разоблачающего рода аргументов. Условия социального выбора интерпретируются как условия убеждений о том, что следует делать, а не предпочтений, благодаря чему обнаруживается, что они имеют эпистемологическое значение. Исследуются также тесные связи между когнитивными концепциями сапоопределения или автономии и демократии, а также, как естественное продолжение идеи этики, развивается идея полностью структурированной социальной теории, считающей внутри- и межличностные отношения одинаково базовыми. Наконец, показывается влияние вопросов о внутриличностной структуре из теории принятия решений, на вопросы справедливого распределения.

Эта книга организована подобно спирали или даже нескольких переплетающихся спиралей: некоторые идеи и аналогии рассматриваются несколько раз и с разных сторон, и нити аргументов, составляющие книгу, переплетаются воедино. Так, аналогия между внутри- и межличностными отношениями очевидна уже на основе приведенного выше очерка. Регулярно будут всплывать различные аристотелевские темы и подчеркиваться естественность рациональности: естественность подходящего содержания, естественность личности и способности к самоинтерпретации и самоопределению. Возможность натурализма, будучи упущена неразборчивым разоблачителем, аннулирует его старания. Более того, стратегия ''разделяй и властвуй'', подразумеваемая в разделении вопросов централизма, субъективизма, когнитивизма, реализма и скептицизма, а также в явном исследовании отношений между ними, имеет важное значение для аргументативной структуры книги. Эти проблемы, полагаю, слишком часто решаются вместе, как если бы за всеми ними стояла какая-то одна большая (хотя, возможно, невыразимая или выразимая лишь метафорически).

Возникающая в результате таких переплетений, сложность структуры, приводит к периодическому повторению, которое я и не пыталась полностью устранить, поскольку надеюсь, что это сделает некоторые главы более самостоятельными и облегчит чтение как читающим только определенные главы, так и читающим с начала до конца. Надеюсь также, что обильные перекрестные ссылки в главах будут полезны некоторым читателям.

Читателям предлагается погружаться в наиболее интересующие их главы. Это, конечно, повышает риск недопонимания, особенно учитывая сильную взаимозависимость глав, но в случае этой книги есть причины поощрять такой подход. Книга эта безусловно междисциплинарная: она ориентирована на аудиторию профессиональных ученых и аспирантов из нескольких разных областей: экономики, права, политической теории и теории принятия решений, а также философии. Профессионалы из каждой из областей найдут здесь хорошо знакомое им и оттого будто бы стандартное, имеющее отношение к их области и могущее представлять непосредственный интерес, а также то, что слишком далеко от их области или слишком технично (впрочем, я сознательно избегала формального изложения технического материала, делая выбор в пользу прозы и доступности). Так, философы наверняка найдут содержание Части I хорошо знакомым, и, возможно, сочтут позволительными и гораздо большие допущения. Хорошо знакомым с теорией принятия решений содержание главы о предпочтениях может показаться излишним. Политические теоретики, возможно, пожелают прочитать длинную заключительную главу --- единственную, посвященную в первую очередь политической теории --- и оценить, какие из предыдущих глав могут представлять для них наибольший интерес. Аналогично, экономисты могут начать с беглого просмотра глав 4, 6 и 12, а юристы --- 10 и 11. Надеюсь, что читатели оценят мои усилия сделать книгу доступной для представителей различных дисциплин и организуют свое ее чтение соответствующим образом.

\part{Разум и ценность}

\chapter{Объективность}

\section{Терминология и разграничения между когнитивизмом, реализмом, централизмом и объективизмом касательно оснований к действию}

В своей яркой и влиятельной книге по этике Мэки использует существование этических разногласий как один из двух основных аргументов в пользу этического скептицизма и субъективизма. Успех его, однако, зависит от того, как понимается объективность, и одним из благотворных последствий скептицизма Мэки было стимулирование явного рассмотрения концепций объективности. Особая концепция объективности предлагается в более поздних работах Витгенштейна и в некоторых работах Дэвидсона. В этой главе я обрисую, каким образом взгляд на основания к действию может быть объективистским, а в следующий --- оспорю подрыв так понятого объективизма разногласиями. В следующих трех главах я разовью аргументы в пользу объективизма в отношении оснований к действию.

Здесь следует сделать несколько предостережений: во-первых, я буду отличать вопрос объективности от прочих, пересекающихся с ним, и попытаюсь защитить объективизм от аргументов от разногласий, подразумевая интерес в нем. Подробно аргументировать такой объективизм я буду уже в последующих главах, особенно в тех, что посвящены предпочтениям, интерпретации и субъективизму. Во-вторых, предполагается некоторые знакомство читателя со взглядами Дэвидсона и позднего Витгенштейна --- в этой книге вы не найдете их систематического изложения и защиты. Думаю, эти взгляды достаточно распространены и влиятельны, чтобы следующее из них понимание оснований к действиям представляло интерес даже будучи, разумеется, подвержено тем же что и они возражениям. В этой и следующей главе я надеюсь показать отсутствие трудностей у этого взгляда с разногласиями. Итак, большая часть разъяснений касательно взглядов Витгенштейна и Дэвидсона здесь будет предназначена для напоминания, выделения и интерпретации в рамках решения поставленной задачи, а не для посвящения в них новичка или переубеждения их противника: мои аргументы касаются скорее связи между их взглядами и защищаемым мной пониманием оснований к действиям. В-третьих, я все же не претендую на приписывание моей позиции Витгенштейну и Дэвидсону.

В отношении объективности ценностей, управляющих действиями, часто рассматриваются по крайней мере два вопроса. Первый --- о том, что я назову когнитивизмом --- касается наличия у суждений о том, что следует сделать, ассерторической или же императивной силы, т.е. выражения ими истинности / ложности и, значит, возможности знания, или же другого рода отношений --- предпочтений, например. Второй --- о том, что я назову научным реализмом --- касается первенства оценочных концепций в смысле их использования в идеальных, научных причинных теориях мира, включая наше использование оценочных понятий. Если причинное объяснение мира, включая использование нами определенных концепций, само по себе зависит от использования таких концепций, то гарантирован реалистский взгляд на предмет утверждений использующих эти концепции; если же такие концепции не фигурируют в идеальных теориях мира, то реализм не оправдан.

Проблема объективности в рассматриваемом мной смысле окажется пересекающейся с проблемами когнитивизма и реализма. Часто объективность поспешно отрицают вслед за когнитивизмом или реализмом. Можно сохранять объективность утверждений о том, что следует сделать, признавая их концептуальную связь с действием, выражение ими предпочтений и отсутствие в них первичных концепций. Эти замечания будут подробно разобраны и подтверждены в последующих главах после решения вопроса объективности.

Возможность объективности тесно связана с концептуальным приоритетом и независимостью. Концептуальная оценка X --- отчет о подразумеваемом, понимаемом и намереваемом при речи или мысли об X. Если об X нельзя дать такой отчет в терминах Y, то X концептуально независим Y, и если при этом подобный отчет об Y может быть дан в терминах X, то X концептуально первичней Y. Вопросы концептуального приоритета и независимости тесно связаны с вопросами о том, каковы вещи и их конститутивные друг с другом отношения, и их следует различать с вопросами приоритета эпистемического и того как мы можем сказать, каковы вещи. (Например, говоря, что знания не находятся в голове, Патнэм отрицает концептуальную независимость значений от мира. Это утверждение о том, что мы имеем в виду, понимаем и о чем намереваемся говорить или думать, когда говорим или думаем о том, что имеем в виду. Это также и конститутивное утверждение о разумеемом под пониманием нечта ''водой''. Оно не о происхождении знания об имеемом в виду или способе определения значений.) Отрицать претензии на концептуальную независимость --- значит отрицать возможность правильного понимания постулируемого с использованием одной или нескольких концепций без понимания постулируемого с использованием другой или нескольких других. (Например, можно говорить о существовании конкретного значения как о нахождении в определенных отношениях с миром.) Чтобы понимать концепцию, не обязательно совершенно и неизменно верно применять ее --- достаточно уметь верно применять ее в некоторых нужных случаях. Отрицание независимости, значит, не предполагает незначительность правильного использования концепций-зависимостей для понимания зависимой концепции. И отрицание приоритетности не подразумевает утверждение обратной приоритетности --- возможна взаимозависимость. В большей части этой книги, за исключением указания обратного (например, в Разделе 3 Главы 15), говоря о независимости и приоритете, я имею в виду не эпистемические, а концептуальные и конститутивные отношения.

Характерной для многих философских подходов к этике является понимание общих --- таких как право и должное --- концепций концептуально первичными и независимыми от конкретных --- таких как справедливое и доброе. Общие концепции, согласно ним, несут фундаментальное значение, которое может быть связано либо с ассерторической, либо с императивной силой, что придает конкретным концепциям статус обосновывающих схватыванием правильного или его доказательством. Подходы, считающие общие понятия концептуально первичными и независимыми, я буду называть централистскими. Пример когнитивистской версии централизма можно найти в работах Росса, пример нон-когнитивистской --- в творчестве Хэйра.

Нон-централизм касательно оснований к действиям, значит, отрицает концептуальную первичность и независимость таких общих понятий как право и должное от таких конкретных обосновывающих концепций как справедливость и недоброжелательность. Он может начать с идентификации таких дискретных конкретных ценностей как справедливость и доброта, и дать отчет о взаимозависимости между общими и конкретными обосновывающими концепциями. Пример нон-централизма --- когерентизм. Согласно когерентизму, должность совершения действия значит его согласованность с теорией, дающей лучшее среди альтернатив объяснение отношений между конкретными ценностями. Такой взгляд должен, значит, по крайней мере неявно, указывать характеризовать эту теорию отношений между конкретными ценностями.

Стоит отметить, что нон-централистская концепция согласованности --- не форма редукционизма, хотя и утверждает супервентность оценочных концепций на неоценочных, т.е. одинаковость в терминах значимого неоценочного влечет одинаковость в терминах оценочного. Когерентизм, таким образом, дает естественное объяснение совместимости супервентности и нередуцируемости, вызывая у некоторых недоумение. Согласно когеретнизму, должность совершения определенного действие не значит его одобрение какой-либо данной теорией или указание наиболее когерентной теории --- нет --- она значит лишь его одобрение наиболее когеретной теорией, которая должна быть открыта уже апостериорно. Супервентность же общих понятий на тех, приложения которых составляют предмет рассматриваемой теории --- сопутствующий фактор их теоретического статуса. Лучшая теория (или теории, в случае ничьей) неизменна от того, о чем она, иначе не будет выполнять объяснительную функцию. Если общие понятия выражают требования лучшей теорией отношений между конкретными ценностями в данных обстоятельствах, то общие понятия ожидаемо супервентны на конкретных оценочных и необходимых для описания обстоятельств применения. Если же конкретные оценочные концепции супервентны на (или, возможно, редуцируемы к) неоценочным, то общие концепции ожидаемо супервентны на неоценочных. Супервентность, таким образом --- признак теоретического, объяснительного. Но поскольку обращение к теории само по себе ничего не дает для определения лучшей теории, требование супервентности не нуждается в редукционистском (или нон-когнитивистском) обосновании. Какая теория лучшая --- уже следующий вопрос, и ответ на него не дается обращением к лучшей теории, какой бы она ни была.

Я не собираюсь более полно разрабатывать понятие супервентности до главы о скептицизме. Сейчас же, для понимания альтернативности когерентизма централизму, полезно объяснить один момент более формально. Когерентист считает, что говоря о правильности конкретной альтернативы, мы подразумеваем существование некой лучшей теории о конкретных ценностях, которая применима к этим альтернативам и благоприятствует одной из них. Такое утверждение следует отличать от более сильного редукционистского о существовании лучшей теории о конкретных применимых к альтернативам ценностях и правильности альтернативы как подразумеваемой благоприятствованием ей этой теории. Первое утверждение сохраняет супервентность и не влечет редукционизм, в отличие от второго, отличая ''необходимо существует такая теория, что ...'' и ''существует такая теория, что необходимо ...''. В каждом возможном мире может существовать некоторая выполняющая требуемую работу теория, но не существует выполняющей эту работу во всех возможных мирах (за исключением, разве что, неинтересно редуцирующей дизъюнктивной).

Рассмотрим образец централистского подхода к такому обосновывающему концепту как справедливое через чисто описательный 'справедливое' плюс общий оценочный O, представляющий стандартный деонтический оператор указания на должное. Кавычки здесь значат воздержание от обыкновенно содержащейся в понятии оценки, сохраняя, таким образом, исключительно описательный его компонент. Общий же оценочный компонент присоединяется к альтернативам в силу их чисто описательных, т.е. неоценочных, характеристик. Согласно некоторым нон-когнитивистским централистским взглядам, общий оценочный компонент скорее императивен чем ассерторичен. Но централистским подход делает не нон-когнитивизм как таковой, а концептуальная независимость и первичность, которые он придает общему оценочному компоненту в отношении конкретных обосновывающих: справедливое в приведенном примере понимается как общее O плюс чисто описательный компонент. Когерентизм же предполагает взаимозависимость общих и частных концепций.

Централизм привлекателен в том числе облегчением понимания критики конкретных этических стандартов. Непарадоксальный статус фразы ''Возможно, это было грубо / нецеломудренно / непатриотично, но не было неправильно'' легко объяснить концептуальной независимостью правильного и неправильного от конкретных ценностей и возможности свободного и внятного их связывания с любыми чисто описательными характеристиками покуда мы делаем это последовательно. Таким образом, одна из задач нон-централизма --- формулирование своего описания концептуальных отношений между конкретными обосновывающими и общими концепциями, которое позволило бы избежать парадоксальности подобных приведенному утверждений. Но проблема эта может быть решена когерентистским пониманием правильного и неправильного как как отражающих статус в рамках лучшей теорией взаимосвязей между различными ценностями. Несоответствие альтернативы некоторой конкретной ценности не противоречит этой альтернативы соответствию лучшей теории взаимосвязей между всеми значимыми конкретными ценностями.

Ключевое различие между централизмом и нон-централизмом, значит, в первичности и независимости общего: централист считает общие концепции первичными к и независимыми от конкретных обосновывающих и объясняет обосновывающий статус последних в терминах общих понятий, тогда как нон-централист отвергает такое притязание и должен дать отчет о концептуальных отношениях между знакомыми конкретными и общими концепциями. В частности, в случае когерентизма, рассматриваемые отношения считаются таковыми предмета и теории, и в главе о когерентности я рассматриваю, как именно следует характеризовать такие отношения. Но отрицание первичности общих понятий не следует первичность конкретных, так же как из независимости данных от теории не следует независимость теории от данных: данные могут быть теоретическими, рефлексивный эквилибриум может быть широким, частные и общие понятия могут быть взаимозависимы.

Итак, я говорила об отличии проблем когнитивизма и реализма от проблем объективизма и о связи последнего с проблемой централизма. Но как соотносятся централизм и объективизм? Оставшаяся часть этой главы и следующие четыре будут посвящены этому вопросу.

Есть нечто правдоподобное в наличии концептуальная связи между общими оценочными концепциями и предпочтениями. Нон-когнитивистская версия этого взгляда утверждает, что общие оценочные концепции служат для выражения предпочтений. (Что именно в нем правдоподобно я разбираю в Разделе 3 этой главы и в главах посвященных субъективизму и когнитивизму.) Нон-когнитивизм не влечет централизм, но если объединить функцию общих понятий как выражающих предпочтения с централизмом (как это обычно и делается), можно получить субъективизм в отношении оснований к действиям: предпочтения первичней конкретных ценностей и независимы от них и в некотором роде их определяют. Объективизм в отношении оснований к действиям --- это отрицание независимости, первичности, или определяющего статуса предпочтений по отношению к ценностям. Таким образом, когерентизм --- версия объективизма, в отличие от субъективизма утверждающая взаимосвязанность предпочтения и ценностей. (Заметьте: можно отрицать субъективизм не отрицая нон-когнитивизма, если отрицать централистское утверждение концептуальной независимости предпочтений от ценностей; этот вид объективизма я рассматриваю в Разделе 3.)

Обратные же утверждения о первичности, независимости и детерминированности --- первичность, независимость и определяющий статус ценностей по отношению к предпочтениям --- я буду называть платонизмом. Платонизм, строго говоря --- версия объективизма, но, конечно, говоря об объективизме я не подразумеваю обязательно платонизм, тем более что среди прочих объективистских позиций он наименее привлекателен. Субъективисты, считающие платонистов своим основным оппонентом, зря тратят время и недооценивают силу своей же позиции. Доктрина взаимозависимости, с которой я связываю объективизм, не рассматривает предпочтения и ценности как независимые и потому не порождает общей для субъективизма и платонизма проблемы. Этот шаг можно считать частью аналогичного более основательного шага касательно разума и мира. Я выступаю против субъективизма и утверждаю, что, по крайней мере в этике, платонизм нежизнеспособен, так что аргументы в пользу взаимозависимости подкрепляются еще и аргументами против субъективизма.

\section{Параллели: централизм касательно цвета, права и логики и некоторые предварительные сомнения}

Вопрос централизма не так уж странен и чужд, как может показаться на первый взгляд. Его привлекательность в отношении этических, юридических или логических оснований, возможно, сильнее, чем в отношении цвета, но даже централизм касательно оснований вызывает некоторые вопросы. В этом разделе я сделаю несколько предварительных комментариев по мере проведения параллелей, продемонстрировав некоторые из таких вопросов и задав контекст для последующих разделов. В следующем разделе я начну рассматривать применение аргументов Витгенштейна и Дэвидсона против определенных централистских позиций, возможность аналогичного противостояния централизму касательно оснований к действиям и ее связь с объективизмом.

Централизм в отношении цвета значит первичность понятия цвета понятиям конкретных цветов и ее от них независимость. Эти конкретные понятия --- такие как красный и зеленый --- говорят о цветах посредством наличия в их значении соответствующего компонента, первичного к различному в них и независимого от этого. Я буду использовать кавычки с понятиями конкретных цветов для обозначения соответствующего ''не-цветного'' описания (в терминах длин волн, например). То есть централизм говорит о красном как о 'красном' плюс наличие воспринимаемого людьми при нормальных условиях как цветное.

Чем мотивирован централизм касательно цвета? Кавычки предполагают своего рода отстраненное размышление о способности непосредственно воспринимать цвет, здравое понимание вторичного статуса понятий цветов и их концептуальной связи с опытом. Возможно, говоря о цветах, мы имеем в виду именно отражающее их второстепенный статус, но централистское понимание цветов отражает этот их статус особым образом: оно не оставляет в стороне вклад аппарата восприятия, а, скорее, 'округляет' его и изолирует в первичном цветовом компоненте. Поэтому, если целью было отразить второстепенный статус понятий конкретных цветов, то централизм был излишним. Можно было утверждать тенденцию объекта любого цвета вызывать у нормальных людей в нормальных условиях переживания этого цвета. (Аналогично правдоподобно утверждение концептуальных связей между общими оценочными понятиями и предпочтениями.) Но это можно принимать, не принимая независимости характера переживания \textit{цвета как такового} от конкретных цветов. (Аналогично наличие концептуальных связей между общими оценочными понятиями и предпочтениями не влечет централизм.) То есть вторичность цветовых понятий не определяет централизм. Хотя каждому конкретному цвету соответствует положение в спектре длин волн, называние красным, желтым или синим не ограничивается таким соотнесением, и квалитативные различия цветов нельзя выразить через понятие цвета как такового. Понятие цветного зависит от понятий конкретных цветов и отношений между ними, позволяющих их комбинировать для получения других, в том числе невозможных, как, например, пурпурно-желтый, цветов. Именно через такие отношения возможно говорить о незнакомых цветах. (На эту тему будет сказано больше в следующей главе.)

Рассмотрим теперь правовую параллель. Правовой централизм значит первичность общей концепции права конкретным правовым концепциям и связанным с ними принципам и ее от них независимость. Централистское понимание такой юридической концепции как договор, например, будет отражать ее статус как основания для юридического решения с точки зрения общей концепции права, которая первична и независима от конкретных правовых концепций, принципов и соответствующих практик, подобно, возможно, Хартовскому правилу признания. Пусть применение кавычек к конкретному юридическому понятию значит соответствующую юридически нейтральную концепцию или правило поведения, тогда централист понимает договор в терминах удовлетворения правило, устанавливающего условия для 'контракта' плюс подтверждение этого правила правилом признания.

Давайте рассмотрим некоторые возражения правовому централизму. Понимание общей концепции права требует понимания по крайней мере некоторые конкретных правовых концепций и связанных с ними практик, и, значит, признание определенных принципов как обязательных для судов не зависит от признания какого-либо центрального статуса обоснования (например их подтвержденности правилом признания). Будучи подкреплено правилом признания, судебное решение может соответствовать конкретной правовой практике, не подтвержденной основным правилом, и все же остающейся правовой практикой. Таким образом, общая концепция права с точки зрения правила признания сплавлена с конкретными правовыми практиками и не может быть отождествлена с аргументирующим статусом концепций, которые эти практики определяют. Если же аргументирующий компонент расширить и применить также и к обыденной практике, то такое подмешивание можно избежать ценой избыточности.

Рональд Дворкин рассматривает версию правового позитивизма Харта применительно к сложным делам, таким как дела Риггса и Хеннингсона, в которых суд ссылается на принципы для оправдания принятия или применения новой нормы. Дворкин указывает, что суд Риггса, обосновывая новую интерпретацию закона о завещаниях, сослался на принцип, согласно которому никто не может извлечь выгоду из своих проступков, и что суд Хеннингсона сослался на несколько пересекающихся принципов для обоснования нового правила об ответственности производителя автомобилей за их дефекты. Согласно Харту, нормы права валидны от подтвержденности правилом признания, которое само по себе не валидно, а просто принимается, на что Дворкин возражает:

\smallskip

''Принципы Риггса и Хеннингсона не проходят предложенную проверку, потому что их происхождение не в каком-то конкретном решении какого-либо законодательного органа или суда, а в чувстве уместности, с течением времени развивающемся у профессионалов и общества, и власть этих принципов зависит от сохранения такого чувства, ведь если больше не кажется несправедливым позволять людям извлекать выгоду из своих ошибок или справедливым --- возлагать особое бремя на производящие потенциально опасные автомобили олигополии, то принципы эти более не играют роли в новых делах даже не будучи когда-либо отменены. Харт же утверждает, что основное правило может назначить законом не только те правила, что приняты конкретными правовыми институтами те, что установлены обычаями. Многие из наших даже самых древних правовых норм созданы не законодательным органом или судом: впервые появившись в юридических текстах и заключениях, они уже были частью права --- к ним относились как к таковой --- поскольку представляли обыденную практику сообщества или какой-то специализированной его части. (Обычно приводятся примеры правил коммерческой практики вроде регулирующих, какие права возникают в рамках стандартной формы коммерческих векселей.) ... Основное правило, согласно Харту, допускает признание некоторых обычаев законов даже до признания их таковым судами. Но это создает трудности для его теории, которые он даже не пытается рассмотреть, поскольку не устанавливает критерии, которые основное правило могло бы для этого использовать. Нельзя использовать единственным таким критерием рассмотрение обществом определенной практики как морально обязательной потому что тогда не будет различия между правовыми и моральными обыденными нормами, тогда как очевидно не все давние традиционные моральные обязательства в обществе диктуются законом. Если же критерий --- убежденность общества в юридической обязательности этой практики, то он обессмысливает основное правило по крайней мере для этого класса правовых норм. Появление основного правила, утверждает Харт, знаменует переход от примитивного общества в правовому, ведь дает возможность определять социальные нормы права иначе чем взвешивая согласие с ними. Но если основное правило утверждает юридическую обязательность любых считаемых так обязательными правил, то ничего оно не определяет и становится в этих случаях не-правилом, ведь и в каждом первобытном обществе есть такое вторичное правило признания, согласно которому обязательно все что таковым считается.''

\smallskip

Наконец, давайте рассмотрим параллель с логикой. Логически обоснованный вывод сохраняет истинность при переходе от посылок к заключению: если посылки истинны и переход верен, то истинно и заключение. Но есть разные валидные переходы: например, делая из \textit{p} и \textit{если p, то q} переход к \textit{q}, мы используем \textit{modus ponens}. Логический централизм предполагает первичность общей концепции сохранения истины, ее независимость от конкретных логических характеристик переходов вроде ''\textit{... --- пример modus ponens}''. Рассмотрим понимание характеристики ''\textit{является примером modus ponens}'' как подтверждающей переходы исходя из концептуально первичной к и независимой от конкретных характеристик переходов общей концепции сохранения истинности. Такая общая концепция может быть принята, например, конвенционалистами как выражения одобрения в обществе определенного нейтрально описанного способа совершения перехода. Пусть закавычевание конкретной логической характеристики вывода ('\textit{modus ponens}', например) дает семантически нейтральную характеристику в чисто синтаксических терминах. Тогда централизм значит понимание перехода modus ponens с точки зрения удовлетворения этого перехода правилу '\textit{modus ponens}' и принятия всего что удовлетворяет этому правилу согласно центральному понятию сохранения истинности.

Взгляд Витгенштейна о значениях как определяемых практиками, будучи применен к логическим константам и математическим концепциям, приводит к описываемому Криспином Райтом как доктрина первичности логики и математики по отношению к истинности. Согласно этой доктрине, ''бессодержательна идея чего-то как в действительности следствия некоторого набора утверждений не только лишь в смысле валидности необходимых переходов'', ''за нашими процедурами вывода нет ничего что было бы ''правильной'' концепцией такого вывода''. Так, Витгенштейн пишет:

\smallskip

''Шаги, не подвергаемые сомнению --- логические умозаключения. Но сомнению они не подвергаются не потому что ''заведомо соответствуют истине'' --- или что-то в этом роде --- нет, просто именно это и называется говорением, думанием, выводом, аргументированием. О каком-либо соответствии между сказанным и действительностью здесь речи нет, скорее логика предшествует любому такому соответствию так же как установление метода измерения предшествует правильности или неправильности определения длины.

--- Но разве нет истины, соответствующей логическому умозаключению? Разве не \textit{истинно} что одно следует из другого?

--- ''Истинно что одно следует из другого'' значит что одно следует из другого.

... 'считать правильно' ... значит считать \textit{так}.''

\smallskip

Хотя доктрина первичности логики перед истинностью влечет отрицание логического централизма, Витгенштейн не ограничивает определяющую силу практик влиянием только на логические и математические понятия. Определяемость содержания практиками не допускает возможности практики быть оспоренной исключительно путем применения той самой концепции, содержание которой она определяет: ''Если вы измеряете размер стола линейкой, измеряете ли вы также размер линейки? Если измеряете линейку, то не можете в то же время измерять и стол''. Здесь не возникает вопроса о соответствии нашей концептуальной схемы фактам, понимаемым как включающие не только объекты, о которых говорим, но и все что может быть о них верно сказано. Потому что жизнеспособная теория объективной истины и соответствия ей не подразумевает самодавлеющесть значений, их автономность и независимость от контекстов вообще и наших практик в частности.

\smallskip

''Слова ''правильный'' и ''неправильный'' используются при указании действовать в соответствии с правилом. Слово ''правильно'' заставляет ученика следовать далее, слово ''неправильно'' сдерживает его. Можно ли объяснить эти правила ученику, сказав: ''Это согласуется с правилом, а это --- нет''. Ну да, если у него есть понятие такой согласованности. Но что, если оно еще не сформировано? (Дело в том, как как он реагирует на слово ''согласуется''.)

Никто не учится подчиняться правилу через научение пользованию словом ''согласуется''.

Скорее наоборот --- значение слова ''согласуется'' схватывается через научение следовать правилу.''

\smallskip

С этой обобщенной нон-централистской точки зрения нет смысла предполагать существование какого-то независимого от всякого конкретного случая критерия --- некого стандарта значений, оторванного от всех наших встроенных в контекст практик, которому значения эти стремятся соответствовать.

\section{Применение аргументов Витгенштейна и Дэвидсона против централизма; пересечение проблем централизма и когнитивизма}

В оставшейся части этой и нескольких последующих глав я попытаюсь показать, как концепция оснований к действию как объективных связана с решением вопроса о централизме в витгенштейновском духе --- в пользу нон-централизма касательно оснований к действию. В ходе этого я выявлю следствия утверждения возможности плодотворного рассмотрения способа направления языкового высказывания смыслом как частного случая направления действия основаниями к нему. Таким образом, я попытаюсь применить соображения Витгенштейна о способе определения значений и их нормативном отношении к языковому действию --- к способу определения оснований вообще и их нормативному отношению к действию. Я также проведу аналогию, предложенную Дэвидсоном, между проблемой взаимозависимости убеждений и значений при интерпретации языкового действия проблемой взаимозависимости убеждений и желаний при интерпретации действия вообще. И я укажу связи между проблемами Витгенштейна и Дэвидсона: концепция объективности представлена обобщенным принципом неблагосклонной трактовки, и я рассмотрю проблему несогласия --- проблему для витгенштейновского взгляда на основания к действию --- через интерпретацию необходимости согласия в форме жизни с точки зрения необходимой роли благосклонности.

То, каким образом соображения о значении и его отношении к убеждению могут быть применены к вопросам о желании и его отношении к убеждению, выявляется путем рассмотрения совместимости нон-когнитивизма с централизмом. Понимание утверждений о должном действии как о концептуально влекущих императивы, выражающих отношение к действию через скорее предпочтения чем убеждения --- не объясняет отношение всего этого концептуального комплекса к конкретным ценностям. Как отмечалось выше, поскольку сочетание централизма и нон-когнитивизма дает субъективизм, субъективизм можно отвергнуть через его централистский, а не нон-когнитивистский компонент (хотя в конечном счете я выступлю и против нон-когнитивизма, но причины этого не излагаются до Главы 9). Давайте рассмотрим, как можно развить это утверждение.

Например, Майкл Даммит предположил существование альтернативы интерпретации Витгенштейна как отрицателя ''идеи, общей для большинства философов, писавших о значении --- наличия некой одной ключевой концепции теории значений'', ''что есть некая одна особенность предложения, которая определяет его значение''. Альтернатива --- понимание Витгенштейна не как отвергающего идею существования ключевого понятия, а как исходящего из ''отсутствия ключевого понятия на стороне обоснования высказывания, в отличие от понятий истинности, проверки, подтверждения и т.д. --- оно находится скорее на стороне его последствий. Т.е. знание значения предложения --- это знание конвенциональных последствий его произнесения, языковой и неязыковой реакции на него со стороны слушателей а также ответственности говорящего его''.

Аналогично понимание обосновывающих концепций как единых от своих последствий --- например, от следования из них императивов, при определенных условиях может стать альтернативой отказу от централизма в отношении оснований к действиям вообще. Конечно, о влекомости императивов обосновывающими концепциями можно говорить централистски, но не обязательно, потому что нон-централисту не нужно отрицать эту влекомость. Но тогда сама концепция действия, наделяющая выражения императивной силой, не может считаться концептуально первичной к и независимой от конкретных обосновывающих концепций.

Но это не диковинная точка зрения. Конечно, в чем-то она параллельна известным в философской литературе взглядам, и знакомые с такой литературой смогут лучше эти параллели увидеть, если сначала в общих чертах изложить мотивацию такой точки зрения в несколько этапов, которые и представлены в последующих главах. Сначала я дам чрезвычайно абстрактный и программный трехэтапный набросок стратегии, который затем будет частично проиллюстрирован сразу и еще много раз в последующих главах. Тем же, кто совсем не знаком с соответствующей литературой, эскиз этот может быть бесполезен, и им придется дождаться дальнейших иллюстраций. В последующем трехэтапном наброске и последующих главах мы попытаемся наметить несколько наборов взаимосвязанных отношений: между тем что должно быть сделано, действием и предпочтением (проблема когнитивизма); между тем что должно быть сделано и определенными ценностями (проблема централизма); между действием, предпочтением, убеждением, значением, т.е. касающимся разума и естественной и социальной средой (проблемы объективности и зависимости или независимости разума и мира); и между миром и ценностями (проблемы этического реализма и этического скептицизма).

Во-первых, мы отвергнем примитивность и беспроблемность концепции интенционального действия в отличие от, скажем, языкового действия. Как личности мы интерпретируем мир для понимания событий как намеренных действий, выражающих предпочтение точно так же, как интерпретируем его для понимания события как выражающего убеждение действия языкового. Кристофер Пикок пишет: ''Предпочтения, выраженные в поведении, порождают действия, то есть символические события: и эти события, каждое из которых --- пример очень многих типов, не дают никакого ключа к типу, предпочтение к которому выражается этим поведением, пока нет теории интерпретации''. Мы способны интерпретировать символические события как действия, то есть понимать их как выражение убеждений, желаний и намерений с определенным содержанием, но как мы это делаем? Как определяется, к какому типу действия относится данное конкретное событие? Интенциональное действие в целом не более самоинтерпретирующие, чем языковое, и возникает еще одна проблема: как избежать радикальной недоопределенности или неопределенности интерпретации? Некоторые детали того, как эта проблема возникает для предпочтений и действий в контексте теории принятия решений, представлены в главе о предпочтениях.

\end{document}
